\input{../tex-inputs/latex-leseansicht-vorspann}

\section[Arthur Schnitzler an Gustav Schwarzkopf, 28. 7. 1903]{L04439 Arthur Schnitzler an Gustav Schwarzkopf, 28. 7. 1903}
\nopagebreak\mylabel{L04439v}
\rehead{ }\normalsize\beginnumbering\briefempfaengerindex{Schwarzkopf, Gustav@\textsc{Schwarzkopf, Gustav}!zzzSchnitzler, Arthur@\emph{von Arthur Schnitzler}!1903-07-282@{28. 7. 1903}|(be}
\toendnotes[C]{\smallbreak\pagebreak[2]}
\correspDesc{Versand  durch Arthur Schnitzler am 28. 7. 1903 in Hochschneeberg
\newline{}Übermittlung  am 29. 7. 1903 in Hochschneeberg
\newline{}Zustellung  am 30. 7. 1903 in 1/1 Wien 1
\newline{}Erhalt  durch Gustav Schwarzkopf im Zeitraum [29. 7. 1903 – 2. 8. 1903?] in Tiefer Graben 23}\toendnotes[C]{\smallbreak}
\Standort{DLA, A:Schnitzler, HS.1985.1.1897.}
\physDesc{Bildpostkarte, 103 Zeichen
\newline{}Handschrift: Bleistift, deutsche Kurrent
\newline{}Versand: 1) Stempel: »\nobreak{}\oindex{Hochschneeberg@\textbf{Hochschneeberg}|pwk}Hochschnee\textcolor{gray}{berg}, 2\textcolor{gray}{9}. 7. 03, 2–4N\nobreak{}«.   2) Stempel: »\nobreak{}\oindex{I., Innere Stadt@\textbf{I., Innere Stadt}|pwk}\textcolor{gray}{1/1 Wien 1}, \textcolor{gray}{30} 7. {[}03{]}, 8–9½V., Bestellt\nobreak{}«. }\toendnotes[C]{\smallbreak}\pstart{}{\pb}Herrn Guſtav Schwarzkopf\pend{}\pstart{}Wien I.\oindex{Wien@\textbf{Wien}|pw}\pend{}\pstart{}Tiefer Graben 23\oindex{Wien@\textbf{Wien}!I., Innere Stadt@\textbf{I., Innere Stadt}!Tiefer Graben 23@\textbf{Tiefer Graben 23}|pw}.\pend{}{\bigskip}
\pstart
           {\pb}\textcolor{gray}{\textbf{Station Schneeberg\oindex{Bahnhof Hochschneeberg@\textbf{Bahnhof Hochschneeberg}|pw}.}}\hfill \textcolor{gray}{\textbf{Haltestelle Schneebergdörfl\oindex{Station Hengsttal@\textbf{Station Hengsttal}|pw}, 616
                  m.}}\pend
           
\pstart
           \centering{}\textcolor{gray}{\textbf{der Zahnradbahn
                     am Schneeberg\orgindex{Schneebergbahn@Schneebergbahn|pw}, 2075 m.}}\pend
           \vspace{1em}
\pstart
           \raggedleft{}{\pb}28. 7. 903.\pend
           \vspace{0.5em}
\pstart
           Auf Wiederſehen \label{K_L04439-1v}\edtext{Donnerſtag}{\lemma{\textnormal{\emph{Donnerstag}}}\Cendnote{\textnormal{Vgl. A. S.: \emph{Tagebuch}, 30. 7. 1903.}}}\label{K_L04439-1}!\pend
           
\pstart
           Herzliche Grüße!{\\[\baselineskip]}\spacefill\mbox{A.}\pend
           \leftskip=0em{}\selectlanguage{ngerman}\endnumbering\briefempfaengerindex{Schwarzkopf, Gustav@\textsc{Schwarzkopf, Gustav}!zzzSchnitzler, Arthur@\emph{von Arthur Schnitzler}!1903-07-282@{28. 7. 1903}|)be}\mylabel{L04439h}
\begin{anhang}
\end{anhang}\newcommand{\dateiname}{L04439}\newcommand{\titel}{Arthur Schnitzler an Gustav Schwarzkopf, 28. 7. 1903}\newcommand{\editorInnen}{Herausgegeben von Jahnke, SelmaMüller, Martin Anton}\input{../tex-inputs/latex-leseansicht-abspann}
